\documentclass[10pt,a4paper]{amsart}
\usepackage[margin=19mm]{geometry}
\usepackage{amsmath,amssymb,mathtools}
\usepackage[scr=rsfs,cal=euler]{mathalfa}
\usepackage{xfrac}
\usepackage[unicode,hidelinks,bookmarks=false]{hyperref}
\newcommand{\ie}{i.e.\ }
\newcommand{\cf}{cf.\ }
\newcommand{\resp}{respectively\ }
\newcommand{\Subsection}{\subsection}
\providecommand{\nobreakdash}{\nobreak\hbox{-}}
\setlength{\emergencystretch}{2em}
\multlinegap0pt
\usepackage{mathtools}
\usepackage{amssymb}
\usepackage{bm}
\usepackage{enumerate}
\usepackage{xcolor}
\usepackage{float}
\newtheorem{lemma}{Lemma}
\newtheorem{theorem}{Theorem}
\newcommand{\depth}{\operatorname{depth}}
\newcommand{\m}{\mathfrak{m}}
\newcommand{\n}{\mathfrak{n}}
\newcommand{\reg}{\operatorname{reg}}
\title{\large \textbf{Bounds on Hilbert coefficients of Cohen-Macaulay modules having finite projective dimension}}
\author{Samarendra Sahoo}
\date{Source: arXiv:2609.10295v1 (CC BY 4.0)}
\begin{document}
\maketitle

\noindent\emph{Source and licence.} \large \textbf{Bounds on Hilbert coefficients of Cohen-Macaulay modules having finite projective dimension}. Version arXiv:2609.10295v1 (CC BY 4.0), distributed by arXiv under the Creative Commons Attribution 4.0 International licence, http://creativecommons.org/licenses/by/4.0/, as recorded in the arXiv licence metadata for this exact version and shown on the abstract page https://arxiv.org/abs/2609.10295v1. Full licence notice: \texttt{licenses/arxiv-2609.10295-CC-BY-4.0.txt}.\par\smallskip

\noindent\textit{Editorial scope.} Counted unit: the article's theorem main (source0.tex lines 462--469). The article's complete original proof of it is delivered with it (source0.tex lines 470--484, one \texttt{proof} environment of 15 lines). No mathematics is written, repaired or re-derived: the statement and the proof are the article's own lines, in the article's own notation, and no retained line is substituted or re-flowed.  Retained uncounted as the source-local context the proof needs: lines 291--297; lines 354--358; lines 382--384; lines 407--413; lines 452--454.  Nothing was omitted from any retained span, so the package carries no omission record.  No retained line contains a citation, so the wrapper carries no bibliography and no bibliography entry.  No retained line contains a figure environment, an image-inclusion command or drawing code, so no image is delivered and the package declares no asset dependency.  Editorial formatting only, in addition: an AMS \texttt{amsart} preamble that restates the article's own declarations and macros used by the retained lines, namely the environments \texttt{lemma}, \texttt{theorem} on separate counters, together with the standard packages those lines need.  The retained environments print in this document as Lemma 1, Theorem 1 in order of appearance, whereas the article prints the same items with the numbering of its own full document.  The complete native source of the article as distributed in the arXiv e-print for this exact version is bundled unchanged as \texttt{sources/209894/source0.tex}.
\begin{lemma}
\label{lowerbound}
Let $(A,\mathfrak{m})$ be a Cohen-Macaulay local ring of dimension $d$ such that $G(A)$ is Cohen-Macaulay. Let $M$ be a Cohen-Macaulay $A$-module of dimension
$r<d$ and finite projective dimension. Set
$c=\operatorname{reg} G(A)$. Then
$e_0(M)\geq\mu(M)+c$ and $e_1(M)\geq\binom{c+1}{2}$.
\end{lemma}

\begin{lemma}\label{kam}
   Let $(A,\mathfrak{m})$ be a Cohen-Macaulay local ring and $M$ be a Cohen-Macaulay $A$-module of dimension $r$. Let $x_1,\ldots ,x_{r-2},x,y$ be $A\oplus M$-superficial sequence. Set  $L=M/(x_1,\ldots ,x_{r-2},x)M$, $N=L/yL$, and  $(B,\n)=(A/(x_1,\ldots ,x_{r-2},x),\m/(x_1,\ldots ,x_{r-2},x))$. Then
   
   $$e_2(M)\leq e_2(N) + \sum_{i\geq 1}i\lambda_B\left(\frac{{\n}^{i+1}L:y}{{\n}^iL}\right) $$ Moreover, if equality holds, then depth $G(M)\geq r-1$.
\end{lemma}

\begin{lemma}\label{implemma}
     Let $(A,\m)$ be a Cohen-Macaulay local ring of dimension $d$ with $G(A)$ Cohen-Macaulay. Let $M$ be a Cohen-Macaulay $A$-module of dimension $r<d$ with finite projective dimension. Let $x_1,\ldots ,x_{r-2},x,y$ be $A\oplus M$-superficial sequence. Set $c=\reg G(A)$, $L=M/(x_1,\ldots ,x_{r-2},x)M$, $N=L/yL$ and  $(B,\n)=(A/(x_1,\ldots ,x_{r-2},x),\m/(x_1,\ldots ,x_{r-2},x))$. If $e_1(M)=\binom{c+1}{2}+s$ then $$ \gamma=\sum_{i\geq 1} \lambda_B \left(\dfrac{\n^{i+1}L:y}{\n^iL }\right)\leq s.$$
\end{lemma}

\begin{theorem}\label{main+1}
    (with the same hypothesis as in Lemma \ref{implemma}) If $e_1(M)=\binom{c+1}{2}+1$ then the following hold. 
     \begin{enumerate}
         \item \label{1+11} If $\gamma=0$ then $G(M)$ is Cohen-Macaulay. Moreover, $e_2(M)=\binom{c+1}{3}$ with the $h$-polynomial $h_M(z) = \mu(M)+2z+z^2+\ldots+ z^{c}.$
         \item \label{1+12} If $\gamma=1$ then $e_2(M)\leq \binom{c+1}{3}+(c-1).$ Further, if equality holds, then $\depth G(M)= r-1$ and $h_M(z)=\mu(M)+z+\ldots+z^{c-2}+2z^{c}$.
     \end{enumerate} 
\end{theorem}

\begin{lemma}\label{length}
Let $(A,\m,k)$ be a local ring and let $E$ be a finitely generated $A$-module. Then $$\lambda_A({\m E}/{\m^2 E})=1\implies \lambda_A({\m^2 E}/{\m^{3} E})\leq 1.$$
\end{lemma}

\begin{theorem}\label{main}
    (with the same hypothesis as in Lemma \ref{implemma}) If $e_1(M)=\binom{c+1}{2}+2$ then the following hold. 
     \begin{enumerate}
         \item \label{2+11} If $\gamma=0$ then $G(M)$ is Cohen-Macaulay. Furthermore,   $e_2(M)=\binom{c+1}{3}$ with the $h$-polynomial $h_M(z)=\mu(M)+3z+z^2+\ldots +z^c$. 
         \item \label{2+12}  If $\gamma=1$ then $e_2(M)\leq \binom{c+1}{3}+(c-1).$ Further, if equality holds, then $\depth G(M)= r-1$ and $h_M(z)=\mu(M)+2z+z^2+\cdots+z^{c-2}+2z^c$.
          \item \label{2+13} If $\gamma=2$ then $e_2(M)\leq \binom{c+1}{3}+2(c-1).$ Further, if equality holds, then $\depth G(M)= r-1$ and $h_M(z)= \mu(M)+z+z^2+\cdots+z^{c-2}+2z^c$.
     \end{enumerate} 
\end{theorem}

\begin{proof}
   (\ref{2+11}) If $\gamma=0$, then, by the proof of Theorem~\ref{main+1}, $G(M)$ is Cohen-Macaulay. It follows that $h_M(z)=h_N(z)$ and $ e_1(M)=e_1(N)=\binom{c+1}{2}+2$. Therefore, by the Lemma \ref{lowerbound}, there are exactly two possible $h$-polynomials for $M$, namely $h_M(z)=\mu(M)+3z+z^2+\cdots+z^c$ or $h_M(z)=\mu(M)+z+2z^2+z^3+\cdots+z^c$. Replacing $E$ by $N$ in Lemma \ref{length}, we obtain that the second case is not possible. Therefore, $h_M(z)=\mu(M)+3z+z^2+\cdots+z^c$ and $ e_2(M)=\binom{c+1}{3}$.

\noindent
   (\ref{2+12})  If $\gamma=1$, then $G(L)$ is not Cohen-Macaulay. Therefore, $G(M)$ is not Cohen-Macaulay. By Lemma \ref{lowerbound} and the proof of Theorem \ref{main+1}, we have $ e_1(N)=\binom{c+1}{2}+1$ and $h_N(z)=\mu(M)+2z+z^2+\ldots+ z^{c}.$ Hence $ e_2(N)=\binom{c+1}{3}$. By the same argument as in the proof of Theorem \ref{main+1}, we obtain $  e_2(M)\leq \binom{c+1}{3}+(c-1)$. Further, if equality holds, then $\depth G(M)= r-1$ and $h_M(z)=h_L(z)=h_N(z)-(1-z)z^{c-1}
=\mu(M)+2z+z^2+\cdots+z^{c-2}+2z^c$.

\noindent
(\ref{2+13}) If $\gamma =2$, then $G(L)$ is not Cohen-Macaulay. Therefore, $G(M)$ is not Cohen-Macaulay. By Lemma \ref{lowerbound} and the proof of Theorem \ref{main+1}, we have $ e_1(N)=\binom{c+1}{2}$ and $h_N(z)=\mu(M)+z+z^2+\ldots+ z^{c}.$ Hence $ e_2(N)=\binom{c+1}{3}$. 

Since $\gamma=2$, we have two possibilities,  for some $i=\alpha$, $\lambda_B \left(\dfrac{\n^{i+1}L:y}{\n^iL}\right)=2.$ or there is $\beta < \alpha$, such that  $\lambda_B \left(\dfrac{\n^{i+1}L:y}{\n^iL}\right)=1$ for $i=\beta ,\alpha$. In both the case,  we have $\displaystyle  \sum_{i\geq 1}i\lambda_B\left(\frac{{\n}^{i+1}L:y}{{\n}^iL}\right)\leq 2\alpha$. Note that, in the second case, the inequality is strict.

Now, By the same argument as in the proof of Theorem \ref{main+1}, we obtain $ e_2\leq \binom{c+1}{3}+2(c-1)$. Furthermore, if equality holds (i.e. equality in the first case), then by Lemma \ref{kam}, depth $G(M)=r-1$ and $\alpha=c-1.$ This implies that $h_M(z)=h_L(z)=h_N(z)-(1-z)z^{c-1}
=\mu(M)+z+z^2+\cdots+z^{c-2}+2z^c$.
\end{proof}

\end{document}
